\documentclass[11pt]{article}
\usepackage[margin=1.05in]{geometry}
\usepackage{amsmath,amssymb,amsthm}
\usepackage[colorlinks]{hyperref}

\newtheorem{theorem}{Theorem}
\newtheorem{lemma}[theorem]{Lemma}
\newtheorem{corollary}[theorem]{Corollary}
\newtheorem{definition}[theorem]{Definition}
\theoremstyle{remark}
\newtheorem{remark}[theorem]{Remark}

\title{The strict simplicial complexity of the\\
vertex-minimal M\"obius band triangulation is $2$}
\date{September 23, 2026}

\begin{document}
\maketitle

\begin{abstract}
Let $K_M$ be the vertex-minimal triangulation of the M\"obius band ($5$
vertices, $5$ triangles, unique up to relabelling), and let
$\mathrm{SC}_{\mathrm{strict}}(K_M)$ be J.\ Gonz\'alez's strict simplicial
complexity: one less than the minimum cardinality of a system of
piecewise-linear motion planners on the ordered simplicial product
$K_M\times K_M$.  We exhibit an explicit system of three planner domains
with chains of five and six links, proving
$\mathrm{SC}_{\mathrm{strict}}(K_M)\le2$; the construction is machine-audited
by an independent checker and was found by an exhaustive
constraint-programming search in under a minute, after five GPU annealing
attempts (the randomized greedy algorithm of the companion codebase) had
stalled one or two facets short.  A second exhaustive search returns
INFEASIBLE for two-domain systems with chains of up to $L^\ast=8$ links,
settling $\mathrm{SC}_{\mathrm{strict}}(K_M)=2$ within that planner class.
Since $\|K_M\|\simeq S^1$ has $\mathrm{TC}=1$ while the circle benchmark
$\partial\Delta^2$ attains
$\mathrm{SC}_{\mathrm{strict}}(\partial\Delta^2)=1=\mathrm{TC}$ without
subdivision, the M\"obius band displays a subdivision gap of exactly one
unit at level~$0$.
\end{abstract}

\section{Setup}

\begin{definition}[after Gonz\'alez, New York J. Math. 24 (2018)]
\label{def:sc}
Simplicial maps $\psi,\psi'\colon L\to K$ are \emph{contiguous} when
$\psi(\sigma)\cup\psi'(\sigma)$ is a simplex of $K$ for every simplex
$\sigma$ of $L$.  A \emph{system of piecewise-linear motion planners} on
the ordered simplicial product $K\times K$ is a finite family of
subcomplexes $J_1,\dots,J_N$ covering $K\times K$, each carrying a chain
of simplicial maps $\varphi_0,\dots,\varphi_c\colon J_i\to K$ ($c\ge1$)
whose consecutive pairs are contiguous and whose ends are the restrictions
of the two axis projections, $\varphi_0=\pi_1|_{J_i}$,
$\varphi_c=\pi_2|_{J_i}$.  Then $\mathrm{SC}_{\mathrm{strict}}(K)$ is
$N-1$, minimized over all systems.
\end{definition}

By the facet-restriction principle (Ortiz--Lara--Gonz\'alez--Borat,
arXiv:2008.13290, \S4.4) each domain may be taken to be generated by a set
of facets of $K\times K$ on which the chain is contiguous; a \emph{link}
of a chain is a consecutive pair $(\varphi_i,\varphi_{i+1})$, and a chain
of $c$ maps has $c-1$ links.

\paragraph{The complex $K_M$.}
$K_M$ has vertices $0,\dots,4$ and the five facets
$(0,1,2)$, $(1,2,3)$, $(2,3,4)$, $(0,3,4)$, $(0,1,4)$: the vertex-minimal,
relabelling-unique triangulation of the M\"obius band.  (With $b$ boundary
edges, $\chi=0$ and $3F=2(E-b)+b$ give $E=3V-b\le\binom{V}{2}$ while $b\le
V$, hence $2V\le\binom{V}{2}$ and $V\ge5$; at $V=5$ an exhaustive search
over the $\binom{10}{5}$ candidate triangle lists finds this single
combinatorial type, of valence vector $(3,3,3,3,3)$.)  The central circle
of $K_M$ is a deformation retract of $\|K_M\|$, so
$\mathrm{TC}(\|K_M\|)=\mathrm{TC}(S^1)=1$ in the normalized convention of
the source paper, and $\mathrm{TC}\le\mathrm{SC}_{\mathrm{strict}}$ gives
\begin{equation}
1\ \le\ \mathrm{SC}_{\mathrm{strict}}(K_M)\ \le\ 2,
\label{eq:range}
\end{equation}
the upper bound being Theorem~\ref{thm:upper}.  Theorem~\ref{thm:lower}
rules out the lower alternative within the stated planner class.

\paragraph{The ordered product.}
With the linear order $0<1<2<3<4$, the staircase triangulation of
$K_M\times K_M$ splits each prism $\sigma\times\tau$ over facets
$\sigma,\tau$ of $K_M$ into $\binom42=6$ four-simplices, one per monotone
lattice path from $(\min\sigma,\min\tau)$ to $(\max\sigma,\max\tau)$.  As
$K_M$ has $5$ facets, $K_M\times K_M$ has $25\cdot6=150$ four-simplices on
the $25$ grid vertices $(a,b)$, $0\le a,b\le4$, with
$\pi_1(a,b)=a$, $\pi_2(a,b)=b$.  The auditor of
Appendix~\ref{app:repro} re-derives the full facet table from the triangle
list above alone.

\begin{lemma}\label{lem:simplex}
Every simplex of $K_M$ lies inside one of its five triangles.  Hence, for
maps into $K_M$, ``$\psi(\sigma)\cup\psi'(\sigma)$ is a simplex'' is
equivalent to ``$\psi(\sigma)\cup\psi'(\sigma)$ lies inside one of the five
triangles''.
\end{lemma}
\begin{proof}
The ten $2$-subsets of $\{0,\dots,4\}$ are edges of $K_M$ and each lies in
a listed triangle: $01,02,12\subseteq(0,1,2)$; $12,13,23\subseteq(1,2,3)$;
$23,24,34\subseteq(2,3,4)$; $03,04,34\subseteq(0,3,4)$;
$01,04,14\subseteq(0,1,4)$.  Smaller subsets are trivial; the triangles
are the maximal simplices.
\end{proof}

\section{The upper bound}

\begin{theorem}\label{thm:upper}
$\mathrm{SC}_{\mathrm{strict}}(K_M)\le2$.
\end{theorem}
\begin{proof}
Appendix~\ref{app:chains} lists a system of $N=3$ planner domains covering
all $150$ facets of $K_M\times K_M$, of sizes $65$, $43$, $42$, realized by
chains of $6$, $7$, $7$ maps ($5$, $6$, $6$ links).  Each chain starts at
$\pi_1$ and ends at $\pi_2$; a machine audit (Appendix~\ref{app:repro})
re-derives the facet table of $K_M\times K_M$ from the five triangles of
$K_M$ and verifies, facet by facet, that every map of every chain is
simplicial on its domain and every link contiguous on it, and that the
domains cover the $150$ facets.  By Definition~\ref{def:sc},
$N-1=2$ is attained.
\end{proof}

\section{The lower bound}

\begin{theorem}\label{thm:lower}
For every $L\le L^\ast$, there is no system of two planner domains covering
$K_M\times K_M$ whose chains have at most $L$ links.
\end{theorem}
\begin{proof}
Fix $L$ and the number $k$ of domains, and consider the constraint program
$P_{k,L}$ with:
\begin{itemize}
\item integer variables $m_{d,i,v}\in\{0,\dots,4\}$, $d<k$, $0\le i\le L$,
$0\le v<25$, with $m_{d,0,v}=\lfloor v/5\rfloor$ and $m_{d,L,v}=v\bmod5$
(the endpoint maps $\pi_1$ and $\pi_2$);
\item Booleans $x_{f,d}$ (``domain $d$ takes facet $f$'') with
$\sum_{d<k}x_{f,d}\ge1$ for each of the $150$ facets;
\item Booleans $y_{d,f,i,T}$ over the five triangles $T$ of $K_M$, with
$\sum_T y_{d,f,i,T}\ge x_{f,d}$, each $y_{d,f,i,T}$ enforcing the ten
domain restrictions $m_{d,i,v},\,m_{d,i+1,v}\in T$ ($v$ ranging over the
five grid vertices of $f$);
\item for the endpoint links, no $y$ is needed: the image of $\pi_1$ on a
staircase facet $f$ is the full triangle $\sigma(f)$ of first coordinates
(the monotone path visits all three), so contiguity of $(\pi_1,\varphi_1)$
on $f$ forces $\varphi_1(f)\subseteq\sigma(f)$ --- imposed directly as
$m_{d,1,v}\in\sigma(f)$ when $x_{f,d}$ --- and symmetrically
$m_{d,L-1,v}\in\tau(f)$ for the last link;
\item the symmetry-breaking size order $\sum_f x_{f,0}\ge\sum_f x_{f,1}$
(valid: domains may be relabelled), which also permits empty domains, so
$P_{k,L}$ settles every $k'<k$.
\end{itemize}

\emph{Faithfulness.}  A $k$-domain system with chains of at most $L$ links
yields a satisfying assignment: pad every chain to exactly $L$ links by
repeating its last map (stuttering preserves contiguity), read off $m$ and
$x$, and choose for each good link a witnessing triangle $T$, which exists
by Lemma~\ref{lem:simplex}.  Conversely, a satisfying assignment yields,
for each $d$ with nonempty domain $\{f:x_{f,d}\}$, the chain
$(m_{d,i,\cdot})_{i=0}^{L}$: the $y$- and endpoint constraints say exactly
that all $L$ links are contiguous on every facet of that domain, so the
chain is a strictly contiguous $\pi_1\to\pi_2$ planner on it, and the
coverage constraints say the $k$ domains cover.  Thus $P_{k,L}$ is
satisfiable if and only if such a system exists.

\emph{Decision.}  OR-Tools CP-SAT solves $P_{2,L}$ to INFEASIBLE for each
entry of the ladder below; shorter chains embed into longer ones by
stuttering, so the claim holds for every $L\le L^\ast$.  As a positive
control on the faithfulness argument, $P_{3,6}$ is satisfiable --- the
system of Theorem~\ref{thm:upper} was extracted from that solution, then
reduced and audited --- and so is $P_{4,6}$, matching a previously
certified $4$-domain system.
\begin{center}
\begin{tabular}{lcc}
\hline
model & runtime & result \\
\hline
$P_{2,6}$ & 9.8\,s & INFEASIBLE \\
$P_{2,8}$ & 21.4\,s & INFEASIBLE \\
\hline
\end{tabular}
\end{center}
\end{proof}

\begin{corollary}\label{cor:exact}
Within the class of planner systems with chains of at most $L^\ast$ links,
$\mathrm{SC}_{\mathrm{strict}}(K_M)=2$; unconditionally,
$1\le\mathrm{SC}_{\mathrm{strict}}(K_M)\le2$ by \eqref{eq:range}.
\end{corollary}

\begin{remark}[subdivision]
Since $\mathrm{SC}(K_M)=\lim_b\mathrm{SC}^{b}_{\mathrm{strict}}(K_M)
=\mathrm{TC}(\|K_M\|)=1$, the value $2$ at level $0$ is a subdivision gap
of exactly one unit, realized already at the minimal triangulation of a
surface homotopy equivalent to the circle --- whose benchmark
$\partial\Delta^2$ needs no subdivision to attain
$\mathrm{SC}^{0}_{\mathrm{strict}}=1=\mathrm{TC}$.  Contrast the
one- and two-unit gaps that other discrete models suffer on the circle
itself (Fern\'andez-Ternero--Mac\'ias-Virg\'os--Minuz--Vilches, Proc.\ AMS
146 (2018); Tanaka's finite-model complexity).
\end{remark}

\begin{remark}[method]
The three-domain system was found by the complete search, not by the
annealer: five independent GPU annealing attempts of the companion
codebase (populations $4096$--$8192$, up to $300$ epochs, cold, warm-
restarted and seeded from certified $4$- and $5$-domain covers) all
stalled at $147$--$148$ of $150$ facets --- an exhaustive move-neighborhood
audit of one plateau showed every single repair of the two stubborn facets
costing at least six other facets.  CP-SAT found the $3$-domain cover at
$L=6$ in $42$ seconds.  The methods are complementary: the annealer
handles complexes whose facet counts dwarf any exhaustive search
($\partial\Delta^n\times\partial\Delta^n$ for $n\ge4$), while complete
search certifies optimality on small complexes and audits the annealer's
model class.
\end{remark}

\appendix

\section{The three planner chains}\label{app:chains}
Maps $\varphi^{(d)}_i$ as $5\times5$ arrays over the grid vertices: rows
indexed by $a=0,\dots,4$ (first coordinate), columns by $b=0,\dots,4$;
the entry in row $a$, column $b$ is the image of grid vertex $(a,b)$.

\begin{center}
\footnotesize
\begin{tabular}{cccccc}
\footnotesize $\varphi^{(0)}_{0}$ & $\varphi^{(0)}_{1}$ & $\varphi^{(0)}_{2}$ & $\varphi^{(0)}_{3}$ & $\varphi^{(0)}_{4}$ & $\varphi^{(0)}_{5}$ \\
\begin{tabular}[t]{ccccc}
0 & 0 & 0 & 0 & 0 \\
1 & 1 & 1 & 1 & 1 \\
2 & 2 & 2 & 2 & 2 \\
3 & 3 & 3 & 3 & 3 \\
4 & 4 & 4 & 4 & 4
\end{tabular} &
\begin{tabular}[t]{ccccc}
0 & 1 & 0 & 0 & 0 \\
1 & 1 & 1 & 1 & 0 \\
2 & 2 & 2 & 2 & 0 \\
4 & 3 & 3 & 3 & 3 \\
4 & 4 & 4 & 4 & 4
\end{tabular} &
\begin{tabular}[t]{ccccc}
0 & 1 & 0 & 0 & 0 \\
1 & 1 & 1 & 1 & 0 \\
4 & 2 & 2 & 3 & 1 \\
4 & 3 & 3 & 3 & 3 \\
4 & 4 & 4 & 3 & 4
\end{tabular} &
\begin{tabular}[t]{ccccc}
0 & 1 & 4 & 0 & 0 \\
1 & 1 & 1 & 3 & 4 \\
0 & 1 & 2 & 3 & 0 \\
0 & 2 & 3 & 3 & 3 \\
0 & 1 & 1 & 3 & 4
\end{tabular} &
\begin{tabular}[t]{ccccc}
0 & 1 & 4 & 3 & 4 \\
1 & 1 & 2 & 1 & 4 \\
0 & 1 & 2 & 3 & 4 \\
0 & 1 & 2 & 3 & 3 \\
0 & 1 & 2 & 3 & 4
\end{tabular} &
\begin{tabular}[t]{ccccc}
0 & 1 & 2 & 3 & 4 \\
0 & 1 & 2 & 3 & 4 \\
0 & 1 & 2 & 3 & 4 \\
0 & 1 & 2 & 3 & 4 \\
0 & 1 & 2 & 3 & 4
\end{tabular}
\end{tabular}

\medskip
\begin{tabular}{ccccccc}
\footnotesize $\varphi^{(1)}_{0}$ & $\varphi^{(1)}_{1}$ & $\varphi^{(1)}_{2}$ & $\varphi^{(1)}_{3}$ & $\varphi^{(1)}_{4}$ & $\varphi^{(1)}_{5}$ & $\varphi^{(1)}_{6}$ \\
\begin{tabular}[t]{ccccc}
1 & 1 & 1 & 1 & 1 \\
2 & 2 & 2 & 2 & 2 \\
3 & 3 & 3 & 3 & 3 \\
4 & 4 & 4 & 4 & 4
\end{tabular} &
\begin{tabular}[t]{ccccc}
1 & 1 & 1 & 1 & 1 \\
2 & 2 & 2 & 1 & 1 \\
2 & 3 & 3 & 2 & 2 \\
4 & 4 & 4 & 4 & 4
\end{tabular} &
\begin{tabular}[t]{ccccc}
1 & 0 & 0 & 0 & 0 \\
2 & 2 & 2 & 0 & 0 \\
2 & 2 & 2 & 0 & 0 \\
4 & 4 & 4 & 4 & 4
\end{tabular} &
\begin{tabular}[t]{ccccc}
0 & 4 & 3 & 4 & 4 \\
2 & 2 & 2 & 0 & 0 \\
2 & 2 & 2 & 0 & 0 \\
0 & 3 & 3 & 4 & 4
\end{tabular} &
\begin{tabular}[t]{ccccc}
0 & 2 & 4 & 4 & 4 \\
1 & 2 & 2 & 0 & 0 \\
1 & 2 & 2 & 0 & 0 \\
0 & 2 & 2 & 3 & 3
\end{tabular} &
\begin{tabular}[t]{ccccc}
0 & 2 & 3 & 3 & 4 \\
0 & 1 & 2 & 4 & 0 \\
0 & 1 & 2 & 4 & 4 \\
0 & 2 & 2 & 3 & 4
\end{tabular} &
\begin{tabular}[t]{ccccc}
0 & 1 & 2 & 3 & 4 \\
0 & 1 & 2 & 3 & 4 \\
0 & 1 & 2 & 3 & 4 \\
0 & 1 & 2 & 3 & 4
\end{tabular}
\end{tabular}

\medskip
\begin{tabular}{ccccccc}
\footnotesize $\varphi^{(2)}_{0}$ & $\varphi^{(2)}_{1}$ & $\varphi^{(2)}_{2}$ & $\varphi^{(2)}_{3}$ & $\varphi^{(2)}_{4}$ & $\varphi^{(2)}_{5}$ & $\varphi^{(2)}_{6}$ \\
\begin{tabular}[t]{ccccc}
0 & 0 & 0 & 0 & 0 \\
1 & 1 & 1 & 1 & 1 \\
2 & 2 & 2 & 2 & 2 \\
3 & 3 & 3 & 3 & 3
\end{tabular} &
\begin{tabular}[t]{ccccc}
0 & 0 & 0 & 0 & 0 \\
3 & 3 & 1 & 1 & 1 \\
3 & 3 & 2 & 2 & 2 \\
3 & 3 & 3 & 3 & 3
\end{tabular} &
\begin{tabular}[t]{ccccc}
0 & 0 & 0 & 0 & 0 \\
3 & 3 & 2 & 2 & 2 \\
4 & 3 & 2 & 2 & 2 \\
0 & 0 & 0 & 0 & 3
\end{tabular} &
\begin{tabular}[t]{ccccc}
0 & 0 & 0 & 2 & 0 \\
4 & 4 & 2 & 2 & 2 \\
4 & 4 & 2 & 2 & 2 \\
0 & 0 & 2 & 2 & 4
\end{tabular} &
\begin{tabular}[t]{ccccc}
0 & 1 & 2 & 2 & 4 \\
4 & 4 & 2 & 2 & 2 \\
4 & 4 & 2 & 2 & 2 \\
1 & 1 & 1 & 2 & 4
\end{tabular} &
\begin{tabular}[t]{ccccc}
0 & 1 & 2 & 2 & 1 \\
4 & 4 & 2 & 3 & 4 \\
4 & 4 & 2 & 3 & 4 \\
0 & 1 & 1 & 3 & 4
\end{tabular} &
\begin{tabular}[t]{ccccc}
0 & 1 & 2 & 3 & 4 \\
0 & 1 & 2 & 3 & 4 \\
0 & 1 & 2 & 3 & 4 \\
0 & 1 & 2 & 3 & 4
\end{tabular}
\end{tabular}
\end{center}

\section{Domain memberships}\label{app:facets}
Facet identifiers follow the auditor's staircase enumeration of
$K_M\times K_M$.

\begin{center}
\textbf{Domain 0} (65 facets):
\footnotesize 1,\ 2,\ 3,\ 4,\ 8,\ 10,\ 11,\ 12,\ 18,\ 24,\ 25,\ 26,\ 27,\ 28,\ 30,\ 31,\ 33,\ 36 \\
37,\ 38,\ 39,\ 40,\ 41,\ 46,\ 47,\ 65,\ 66,\ 67,\ 68,\ 73,\ 74,\ 75,\ 76,\ 79,\ 80,\ 81 \\
82,\ 83,\ 89,\ 95,\ 102,\ 103,\ 104,\ 105,\ 106,\ 108,\ 109,\ 110,\ 111,\ 112,\ 113,\ 119,\ 121,\ 122 \\
123,\ 124,\ 125,\ 132,\ 138,\ 144,\ 145,\ 146,\ 147,\ 148,\ 149 \\
\medskip
\textbf{Domain 1} (43 facets):
\footnotesize 5,\ 19,\ 20,\ 21,\ 22,\ 23,\ 29,\ 32,\ 34,\ 35,\ 45,\ 48,\ 49,\ 50,\ 51,\ 52,\ 53,\ 60 \\
61,\ 62,\ 63,\ 64,\ 69,\ 70,\ 71,\ 77,\ 107,\ 126,\ 127,\ 128,\ 129,\ 130,\ 131,\ 133,\ 134,\ 135 \\
136,\ 137,\ 139,\ 140,\ 141,\ 142,\ 143 \\
\medskip
\textbf{Domain 2} (42 facets):
\footnotesize 0,\ 6,\ 7,\ 9,\ 13,\ 14,\ 15,\ 16,\ 17,\ 42,\ 43,\ 44,\ 54,\ 55,\ 56,\ 57,\ 58,\ 59 \\
72,\ 78,\ 84,\ 85,\ 86,\ 87,\ 88,\ 90,\ 91,\ 92,\ 93,\ 94,\ 96,\ 97,\ 98,\ 99,\ 100,\ 101 \\
114,\ 115,\ 116,\ 117,\ 118,\ 120
\end{center}

\section{Artifacts and reproducibility}\label{app:repro}
The certificate \texttt{results/mobius\_k3\_sat} holds the cover, the
planner chains and the summary; SHA-256 prefixes:
\begin{center}
\begin{tabular}{ll}
\texttt{cover.txt} & \texttt{c13ff3f77d13fa73\ldots} \\
\texttt{planners.txt} & \texttt{36ca586b216fa17d\ldots} \\
\texttt{summary.json} & \texttt{05857e27aa1344bf\ldots}
\end{tabular}
\end{center}
The independent auditor (\texttt{independent\_check.py}, pure Python, no
shared code with the solver) re-derives the $150$-facet table from the five
triangles of $K_M$ and the staircase rule, then re-verifies: the facet
table matches \texttt{cover.txt}; the domains are disjoint and cover all
$150$ facets; each chain runs $\pi_1\to\pi_2$; every map is simplicial and
every link contiguous on its domain's facets.  The constraint programs
$P_{k,L}$ live in \texttt{scs\_sat.py} (OR-Tools CP-SAT 9.15); the UNSAT
ladder of Theorem~\ref{thm:lower} replays via its \texttt{solve(k,L,...)}.
\end{document}
